\subsection{Ordinary local matching in a fixed finite geometric quotient}

\begin{statement}{A.72}{Proved}{Ordinary ghosts retaining a prescribed finite quotient locally}
Let $N\triangleleft\Pi$ be an open subgroup normal in $E_U$, put
$Q=\Pi/N$, and fix the point-$3$ section $s_0$.
Write $A=\Pi^{\rm ab}$ and
\[
 E_{Q,A}=E_U/(N\cap[\Pi,\Pi]),\qquad
 E_Q=E_U/N,\qquad E_A=E_U/[\Pi,\Pi].
\tag{A.72.1}
\]
There is an actual continuous section $\tau:G_\Q\to E_{Q,A}$
with ordinary infinite-support Kummer coordinates $(A_\tau,B_\tau)$
such that its projection to $E_Q$ is $s_0$ and, at every finite
prime $p$, its restriction is geometrically conjugate to the image
of a point $x_p\in U(\Q_p)$.  At the real place its restriction is
the image of the point $3$ up to geometric conjugacy.

The coordinates have the form
\[
 A_\tau=3\prod_{j\geq0}q_j^{n_j},\qquad
 B_\tau=-2\prod_{j\geq0}r_j^{n_j},\qquad
 n_j=N_0(j+4)!,
\tag{A.72.2}
\]
where $N_0$ is a positive integer specified in the proof and all
$q_j,r_j>3$ are distinct.  In particular $\tau$ is neither a rational
point image nor a cusp image in this quotient.
This applies to each of the four actual finite projective Legendre
quotients in A.53, and to any fixed finite collection of such tests
by taking their common finite quotient.
The finite quotient is fixed before the construction; no compatible
construction over all finite geometric quotients is asserted.
\end{statement}
\begin{proof}
We retain the full abelianization and the local conjugacies in the
same arithmetic quotient.  We first prove the extra field-intersection
calculation required to impose a finite cover in the prime recursion.

\emph{A uniform intersection calculation.}
Let $F/\Q$ be a fixed finite Galois extension and set $M=[F:\Q]$.
Choose $C$ divisible by $24$ and $4M$.  For $d_0=3$ or $-2$, put
\[
 K_{d_0}=F\bigl(\mu_{4M},d_0^{1/(4M)}\bigr).
\tag{A.72.3}
\]
These are fixed finite Galois fields.  Choose $f_{d_0}$ with
$K_{d_0}^{\rm ab}\subset\Q(\mu_{f_{d_0}})$, for example the least
such positive integer.  Existence is precisely the finite-extension
case of the Kronecker--Weber theorem
\cite[I, Theorem 4.16, pp.~49--50]{MilneCFT}.
For every $n$ divisible by $24,4M,f_3,f_{-2}$ and every positive
rational $w$ of valuation zero at two and three, set
\[
 d=d_0w^C,\qquad L=\Q(\mu_n,d^{1/n}).
\]
Then
\[
                         (FL)^{\rm ab}=\Q(\mu_n).
\tag{A.72.4}
\]
Here the superscript denotes the fixed field of the closed
commutator subgroup of the finite Galois group over $\Q$.

Indeed, the radical computation in A.44 applies without change to
the $C$th-power factor: its valuation at three or two is still one,
and $8\mid C$.  Thus
$N_L=\operatorname{Gal}(L/\Q(\mu_n))=\langle\sigma\rangle$
has order $n/2$, with $\sigma(d^{1/n})=\zeta_n^2d^{1/n}$.
In $G=\operatorname{Gal}(FL/\Q)$ the subgroup
$N_F=\operatorname{Gal}(FL/F\Q(\mu_n))$ embeds in $N_L$ with
index dividing $M$.  Hence $MN_L\subset N_F$.
Complex conjugation acts by inversion on $N_F$, so
\[
                 2MN_L\subset2N_F\subset[G,G].
\tag{A.72.5}
\]
The fixed field of $2MN_L$ in $L$ is exactly
$\Q(\mu_n,d^{1/(4M)})$, since $4M\mid n$.
The subgroup in (A.72.5) fixes $F$, so its fixed field in $FL$ is
\[
 F\Q(\mu_n,d^{1/(4M)})
   =F\Q(\mu_n,d_0^{1/(4M)})\subset K_{d_0}\Q(\mu_n).
\]
For any finite Galois field $K$ and finite abelian field $T$ over
$\Q$, restriction identifies $\operatorname{Gal}(KT/\Q)$ with the
fibre product over $\operatorname{Gal}(K\cap T/\Q)$.
Every element of $\operatorname{Gal}(K/\Q)$ has a lift to this
fibre product.  Commutators are consequently exactly the subgroup
$[\operatorname{Gal}(K/\Q),\operatorname{Gal}(K/\Q)]\times\{1\}$.
It follows that $(KT)^{\rm ab}=K^{\rm ab}T$.
Apply this with $K=K_{d_0}$ and $T=\Q(\mu_n)$; the conductor
choice proves (A.72.4).

Consequently, for every modulus $m$ divisible by $n$, every unit
class $u\pmod m$ with $u\equiv1\pmod n$ contains infinitely many
primes splitting completely in both $F$ and $L$.
The actual intersection $(FL)\cap\Q(\mu_m)$ is $\Q(\mu_n)$
by (A.72.4), so the pair $(1,u)$ is an element of the Galois
fibre product.  Chebotarev, in the form opened and used in A.44,
supplies these primes and permits any fixed finite set to be omitted.

\emph{The pointed finite cover and its boundary fields.}
The subgroup $N s_0(G_\Q)\subset E_U$ defines a connected finite
etale cover $f:Y\to U$ of degree $q=|Q|$, with a rational point
$y_0$ above $3$.  Geometrically it is the regular $Q$-cover.
Let $\overline Y\to\Pj^1$ be its normalization compactification.
Choose a finite Galois field $F$ over which all points above zero
and one, all their tangent leading coefficients, and the geometric
deck transformations are defined.  Enlarge $F$ to contain
$\mu_q$, and the required roots described next.
At a point above zero of ramification index $e$ write in its
completed local ring
\[
                   t=d z^e(1+z h(z)),
\tag{A.72.6}
\]
and use $1-t$ at a point above one.  Adjoin $e$th roots of $d$
at every such point, and $q$th roots of $3$ and $-2$.
This requires only finitely many algebraic numbers.

Choose an integer $\Delta$ divisible by $6q$ and by the finitely
many primes needed for good models of this cover, its boundary,
and these local parameters.  Outside $\Delta$ the boundary is
finite etale, all indicated units are integral, and all $e$ are
invertible.  The power series $1+zh(z)$ then has a unique $e$th
root of constant term one over each completed local ring.
The change of parameter $z\mapsto z(1+zh(z))^{1/e}$ has unit
linear coefficient.  Hensel's lemma proves the following precise
property: at a prime $p\nmid\Delta$ split in $F$, every sufficiently
small nonzero $e$th-power value of $t$ is reached by every branch
above zero, with all $e$ root choices rational over $\Q_p$.
The analogous statement holds for $1-t$ above one.
With the displayed good integral parameters, positive valuation
already suffices.  Thus the entire fibre of $Y$ splits over the
values used below, not merely one unspecified branch.

Now set $M=[F:\Q]$ and choose $C$ divisible by $24,4M,q$.
Take $N_0$ divisible by $C,\Delta,f_3,f_{-2}$, with the conductors
from (A.72.3).  These choices fix the exponents in (A.72.2).

\emph{The distinguished component and a bound for its good fibres.}
Let $e=\exp(Q^{\rm ab})$, so $e\mid C$.
For coefficients $a=3u^C$, $b=-2v^C$ with positive rational $u,v$
and for $e\mid n$, consider
\[
 F_{a,b,n}:aX^n+bY^n=1,\quad XY\ne0,
       \qquad t=aX^n,
\tag{A.72.7}
\]
and its fibre product with $Y\to U$.
Its largest common geometric quotient is exactly $Q^{\rm ab}$:
the Kummer group is $A/nA$, and the map from $Q$ to any common
quotient is abelian; conversely $e\mid n$ makes $Q^{\rm ab}$ a
quotient of $A/nA$.
Since $a/3$ and $b/(-2)$ are rational $e$th powers, the component
corresponding to the point-$3$ framing of this common quotient is
defined over $\Q$.  Denote it by $W_{a,b,n}$.
More explicitly, its geometric group is
\[
                 Q\mathop{\times}_{Q^{\rm ab}}A/nA.
\tag{A.72.8}
\]
The two common $e$-power coordinates are obtained from
$(a/3)^{1/e}X^{n/e}$ and $(b/(-2))^{1/e}Y^{n/e}$.
Their rational prefactors specify the framing.  More explicitly, let
$Z_Q=Y/[Q,Q]$ be the pointed common abelian cover.  The displayed
coordinates define a $\Q$-morphism from $F_{a,b,n}$ to the pointed
$e$-Kummer cover, and its quotient gives $F_{a,b,n}\to Z_Q$.  Then
$W_{a,b,n}=Y\times_{Z_Q}F_{a,b,n}$.  This fibre product is
geometrically connected because $Z_Q$ is the largest common geometric
quotient.  At $(a,b)=(3,-2)$ it contains $(y_0,1,1)$.
Thus every component and its framing are specified by actual maps.

The degree $d$ of this geometrically connected component is at most
$D=qn^2$.  Its only branch values are zero, one, and infinity.
Riemann--Hurwitz gives $2g\le d+2$, and its boundary has at most
$3d$ geometric points.  At primes outside $\Delta nab$ its smooth
projective completion has good reduction with the same geometric
group and boundary data.  To justify this last comparison, apply
the prime-to-$p$ specialization of the smooth pair
$(\Pj^1,\{0,1,\infty\})$:
\cite[Expos\'e XIII, Corollary 2.9 and \S2.10; proof of
Corollary 2.12, pp.~289--291 in the re-typeset edition]{AgentSGA1}.
Here $p\nmid qn$, so (A.72.8) has order prime to $p$.
The Kummer twists by $p$-adic units extend etale over the base;
the comparison preserves the component and its inertia groups.
Equivalently, the tame local normalization of (A.72.6) supplies
the same smooth compactification after taking the fibre product.

The Hasse--Weil estimate used in A.48 now gives an interior point
whenever
\[
 p>B_Q(n):=(4qn^2+10)^2.
\tag{A.72.9}
\]
Indeed its lower bound is
$p+1-(D+2)\sqrt p-3D$.
At $\sqrt p=4D+10$ this expression is
$12D^2+59D+81>0$, and it increases thereafter.
An interior point lifts to $\Q_p$, with both Kummer coordinates
units, by smooth Hensel lifting.  This proves the required uniform
large-prime assertion for the distinguished component.

\emph{The recursion with local component matching.}
We construct positive-support coefficients
\[
 a_j=3u_j^C>0,\qquad b_j=-2v_j^C<0,
 \quad u_j,v_j\in\Z_{>0},\quad \gcd(a_j,b_j)=1,
\tag{A.72.10}
\]
with $u_jv_j$ prime to six, and put $W_j=W_{a_j,b_j,n_j}$.
For every finite prime, $W_j$ has a point with $X,Y\in\Z_p^\times$.
At every $p\le B_Q(n_j)$ and every $p\mid\Delta$ we additionally
retain a point on this same component with
\[
 t=a_j,\quad X=1,\quad
 Y\equiv1\pmod {p^{v_p(n_{j+1})}},\qquad
 v_p(a_j+b_j-1)\ge V_j(p)+2v_p(n_{j+1})+j+3,
\tag{A.72.11}
\]
where $V_j(p)=\max(v_p(a_j),v_p(b_j))$.
Start with $a_0=3,b_0=-2$ and $(y_0,1,1)$ at every place.

At stage $j$, write $n=n_j$, $m=n_{j+1}$ and let $S$ consist
of the primes at most $B_Q(m)$, those dividing $\Delta$, and all
prime divisors of $a_jb_j$.
For $p\in S$ choose an old unit-coordinate point
$(y_p,\xi_p,\eta_p)\in W_j(\Q_p)$, with
$f(y_p)=a_j\xi_p^n$.
At $p\mid m$ use the point in (A.72.11).
Such a prime is at most $B_Q(n)$: it divides $N_0$ or is at most
$j+5$.  Thus
\[
                   \xi_p\equiv\eta_p\equiv1
                          \pmod {p^{v_p(m)}}\quad(p\mid m).
\tag{A.72.12}
\]

Choose precision integers $L_p$ as large as needed, initially with
\[
 L_p\ge V_j(p)+2v_p(n_{j+2})+j+4.
\tag{A.72.13}
\]
There is also a finite, explicit local requirement on $L_p$:
the new point constructed below must continue the chosen old
component.  We specify why this can always be imposed before
choosing either new prime.  For unit variables $Q_p,R_p$ tending
to $\xi_p,\eta_p$, the value $a_jQ_p^n$ tends to $f(y_p)$.
Since $Y\to U$ is etale at $y_p$, it has a unique local analytic
inverse on a sufficiently small ball about this value.
The normalized polynomial
\[
 Z^m-\frac{1-a_jQ_p^n}{b_jR_p^n}
\tag{A.72.14}
\]
has error at $Z=1$ of valuation at least $L_p-v_p(b_j)$ when
$Q_p\equiv\xi_p$, $R_p\equiv\eta_p\pmod {p^{L_p}}$.
By (A.72.13) it has a root $Z_p$ with
$v_p(Z_p-1)\ge L_p-v_p(b_j)-v_p(m)$.
This estimate tends to infinity with $L_p$.
Writing $k=m/n$, the point
\[
 \bigl(y'_p,Q_p,R_pZ_p^k\bigr),\qquad
                         f(y'_p)=a_jQ_p^n,
\tag{A.72.15}
\]
is therefore arbitrarily close to
$(y_p,\xi_p,\eta_p)$ in the old fibre product.
The chosen geometric component $W_j$ is an open-and-closed
subscheme of that product, so (A.72.15) belongs to it for all
sufficiently large $L_p$.
This is a finite Hensel condition on the fixed old point and
local equations.  Increasing $L_p$ until the corresponding
Jacobian and root-separation inequalities hold is a terminating
choice.  It also gives
$Z_p\equiv1\pmod {p^{v_p(n_{j+2})}}$ by (A.72.13).
This argument includes $p\mid e$: the finite etale common-cover
branches are separated by taking a sufficiently small ball;
no reduction modulo $p$ is used to distinguish those branches.

Put $P=\prod_{p\in S}p^{L_p}$.
It is divisible by $m$.
Using (A.72.4), choose a fresh prime $q_j$ above every prime in
$S$ and above $n_{j+2}$ such that
\[
 q_j\equiv\xi_p\pmod {p^{L_p}}\ (p\in S),\qquad
 q_j\text{ splits in }F\Q(\mu_n,b_j^{1/n}).
\tag{A.72.16}
\]
The CRT class is one modulo $m$ by (A.72.12), hence one modulo
$n$, as required by the proved intersection lemma.
At $q_j$ choose a unit $z$ with
$b_jz^n=1-a_jq_j^n$; complete splitting and simple-root Hensel
lifting give it exactly as in (A.44.16).
Choose a second fresh prime $r_j>q_j$, again by (A.72.4), with
\[
 \begin{split}
 r_j&\equiv\eta_p\pmod {p^{L_p}}\quad(p\in S),\\
 r_j&\equiv z\pmod {q_j},\\
 r_j&\text{ splits in }
        F\Q(\mu_m,(a_jq_j^n)^{1/m}).
 \end{split}
\tag{A.72.17}
\]
The new positive radicand is still $3$ times a $C$th power,
because $C\mid n$.  All hypotheses of the uniform field lemma
hold.  Define $a_{j+1}=a_jq_j^n$ and $b_{j+1}=b_jr_j^n$.

There is an actual morphism of covers
\[
 W_{j+1}\longrightarrow W_j,\qquad
        (y,X,Y)\longmapsto(y,q_jX^{m/n},r_jY^{m/n}).
\tag{A.72.18}
\]
Define $W_{j+1}$ as the component mapping to $W_j$.
It is unique, since the common geometric quotient is the same
$Q^{\rm ab}$ at both exponents; the map on component labels is
the identity under the rational prefactors in (A.72.8).
It is defined over $\Q$.  In terms of the preceding fibre product,
the equality of the common coordinates follows from
$(a_{j+1}/3)^{1/e}=(a_j/3)^{1/e}q_j^{n/e}$
and the analogous identity for $b$ and $r_j$.

For $p\in S$, apply (A.72.14)--(A.72.15) with
$Q_p=q_j,R_p=r_j$.  They give a point
$(y'_p,1,Z_p)\in W_{j+1}(\Q_p)$ with
$f(y'_p)=a_{j+1}$.  The component assertion follows from
(A.72.18) and the already verified membership of (A.72.15).
The coefficient sum error has valuation at least $L_p$, and
$V_{j+1}(p)=V_j(p)$, so (A.72.11) holds at the next stage.

At $p=q_j$, the new Fermat equation has a unit solution with
$X=1$, by the congruence involving $z$ and simple-root Hensel
lifting.  Its value $t=a_jq_j^n$ has positive valuation and is
an $e_0$th power for every zero-boundary ramification index
$e_0$: $e_0\mid q\mid C,n$ and $3^{1/q}\in F\subset\Q_{q_j}$.
Here the last inclusion means a chosen split completion.
Thus the entire $Y$-fibre at this value is rational by (A.72.6).
Every common $Q^{\rm ab}$ branch is available, so one of these
rational lifts lies with the chosen Fermat point in $W_{j+1}$.

At $p=r_j$, take $Y=1$ in the new Fermat equation.
Its unit $X$ exists by the splitting in (A.72.17).
Now $t=1-b_jr_j^n$, and $1-t$ is an $e_1$th power of positive
valuation for every one-boundary index, using $(-2)^{1/q}\in F$.
Again the whole $Y$-fibre splits, so a lift in the designated
component exists.  At every other prime the coefficients and
$\Delta m$ are units and $p>B_Q(m)$, so (A.72.9) supplies a
unit-coordinate point of that component.
This proves all local invariants and completes the induction.

\emph{The full quotient section and all its local conjugacies.}
The products in (A.72.2) converge exactly as in A.44.
The divided products with exponents $n_j/C$ converge as well;
hence $A_\tau/3$ and $B_\tau/(-2)$ are actual $C$th powers.
Their relative Kummer cocycles can therefore be chosen pointwise
in $CA\subset\ker(A\to Q^{\rm ab})$.  Concretely, fix compatible
roots of each selected positive prime and use the literal sums
$\sum_j n_j\kappa(q_j)$ and $\sum_j n_j\kappa(r_j)$.
For every finite coefficient modulus their tails are identically zero,
so these are uniformly convergent sums of continuous cocycles.
The exact fibre-product map in A.51 identifies
\[
 E_{Q,A}=E_Q\mathop{\times}_{E_{Q^{\rm ab}}}E_A.
\tag{A.72.19}
\]
Combining these actual abelian cocycles with $s_0$ in $E_Q$
defines the asserted continuous section $\tau$.

At exponent $n_j$, its finite geometric quotient is (A.72.8),
and the finite cover associated to this section is precisely
$W_j$.  This can be checked on the common coordinates displayed
after (A.72.8): the rational $e$th roots are the same ones used
to form the pointwise $CA$-valued cocycles.
Under an update the additional Kummer cocycles are
$n_j\kappa(q_j)$ and $n_j\kappa(r_j)$, which reduce to zero
at exponent $n_j$.  Thus the sections and their covers are
compatible by the actual maps (A.72.18), including their framing.

For a fixed finite prime $p$, the ordinary limits of $a_j,b_j$
are nonzero numbers $x_p,y_p$ with $x_p+y_p=1$, by (A.72.11).
For each $j$ choose any of the proved unit-coordinate points on
$W_j(\Q_p)$ and let its image be $t_j=a_jX_j^{n_j}$.
The powers $X_j^{n_j}$ tend to one uniformly for
$X_j\in\Z_p^\times$, since $n_j$ eventually annihilates every
finite unit quotient.  Hence $t_j\to x_p$.
For every fixed $i$, (A.72.18) gives lifts of all $t_j$, $j\ge i$,
to $W_i(\Q_p)$.  The image of this finite cover on local points
is closed inside $U(\Q_p)$: use compactness of its projective
compactification and the fact that the limiting value $x_p$ is
away from the three boundary points.  Therefore $x_p$ itself
lifts to every $W_i(\Q_p)$.

A fixed sheet in this pulled-back cover is exactly a geometric
conjugacy between the point section and $\tau$ in the corresponding
finite geometric quotient.  The inverse images of these conjugator
sets in the compact geometric group of $E_{Q,A}$ are nonempty,
closed, and nested.  Their intersection gives one conjugator for
the whole arithmetic section, because the quotients at $n_i$
are cofinal in its abelian factor.  This proves the asserted
joint local matching in (A.72.1), not only matching in its two
projections separately.

At the real place choose positive compatible roots for the divided
products.  Their Kummer cocycles vanish on complex conjugation,
so $\tau$ has the real image of $s_0$.
The fixed valuations at two and three and the two disjoint infinite
supports are exactly those in (A.44.2), with the exponents (A.72.2).
They exclude rational point and cusp images by the explicit
abelian comparison used in A.44.  This finishes the proof.
\end{proof}

The retained certificate \path{code/finite_quotient_local_certificate.py}
checks the cyclic subgroup inclusions in (A.72.5), the genus-bound
polynomial, the coefficient-prime precision inequalities, and the
component-coordinate transition exponents.  Prime-search termination
and the local-cover continuation are proved above; no large numerical
prime stage is claimed as an executed computation.
